\documentclass[11pt,a4paper]{article}
\usepackage[utf8]{inputenc}
\usepackage[vietnamese]{babel}
\usepackage{amsmath,amssymb}
\usepackage[version=4]{mhchem}
\usepackage{geometry}
\geometry{top=1.5cm,bottom=1.2cm,left=1.4cm,right=1.4cm,headheight=28pt,headsep=12pt,footskip=18pt}
\usepackage{graphicx}
\usepackage{xcolor}
\usepackage{tabularx}
\usepackage{array}
\usepackage{tikz}
\usepackage{fancyhdr}
\usepackage{enumitem}
\usepackage{multido}
\usepackage{tcolorbox}
\usepackage{colortbl}

\definecolor{hfRed}{HTML}{D32F2F}
\definecolor{hfGreen}{HTML}{8BC390}
\definecolor{mainGreen}{HTML}{2E7D32}
\definecolor{solBlue}{HTML}{1565C0}
\definecolor{tableGreen}{HTML}{E8F5E9}

\usepackage{ex_test}

\pagestyle{fancy}
\fancyhf{}
\renewcommand{\headrulewidth}{0.75pt}
\renewcommand{\headrule}{\hbox to\headwidth{\color{hfGreen}\leaders\hrule height \headrulewidth\hfill}}
\renewcommand{\footrulewidth}{0.75pt}
\renewcommand{\footrule}{\hbox to\headwidth{\color{hfGreen}\leaders\hrule height \footrulewidth\hfill}}

\fancyhead[L]{\small\itshape\color{hfRed}Tài liệu lưu hành nội bộ}
\fancyhead[R]{\small\color{hfRed}Thầy \textbf{TRẦN VĂN THẠNH} - 0777.470.803}

\fancyfoot[L]{\footnotesize\itshape\color{hfRed}CS1: 6/15 Nguyễn Hoàng, P. Kim Long, TP Huế \quad|\quad CS2: 24 Đặng Thái Thân, TP Huế}
\fancyfoot[R]{\small\textbf{\color{hfRed}Trang \thepage}}

\setlength{\parindent}{0pt}
\setlength{\parskip}{2pt}

\begin{document}

\noindent
\begin{tabularx}{\textwidth}{@{}X r@{}}
\textbf{TRƯỜNG THPT NGUYỄN TRƯỜNG TỘ} & \textbf{ĐỀ KIỂM TRA GIỮA KỲ I} \\
\textbf{TỔ HÓA - TIN} & \textbf{NĂM HỌC 2026 - 2027} \\
\textbf{ĐỀ CHÍNH THỨC} & \textbf{Môn: Hóa học lớp 12} \\
\textit{(Đề có 04 trang)} & \textbf{Thời gian làm bài: 45 phút} \\
& \textbf{Mã đề: 106}
\end{tabularx}

\vspace{2pt}
\noindent\textbf{Họ và tên:} \dotfill \textbf{Số báo danh:} \dots\dots\dots\dots

\vspace{2pt}
\noindent\textit{Cho biết nguyên tử khối của các nguyên tố: H = 1; C = 12; N = 14; O = 16; Na = 23; S = 32; Cu = 64; Ag = 108.}
\vspace{2pt}
\hrule height 0.5pt
\vspace{6pt}

\noindent{\color{mainGreen}\textbf{\Large A. PHẦN TRẮC NGHIỆM}}
\vspace{4pt}

\noindent{\color{mainGreen}\textbf{\large Phần I: Câu trắc nghiệm nhiều phương án lựa chọn.}} \\
\textit{Thí sinh trả lời từ câu 1 đến câu 16. Mỗi câu hỏi thí sinh chỉ chọn một phương án.}
\vspace{4pt}
\begin{ex}
Khi nấu canh cua thấy các mảng ``gạch cua'' nổi lên là do
\choice
{sự đông tụ lipid.}
{phản ứng thủy phân protein.}
{phản ứng màu biuret.}
{sự đông tụ protein.}
\loigiai{}
\end{ex}
\begin{ex}
Công thức cấu tạo của methylamine là
\choice
{\ce{C2H5NH2}.}
{\ce{CH3NH2}.}
{\ce{CH3NHCH3}.}
{\ce{(CH3)3N}.}
\loigiai{}
\end{ex}
\begin{ex}
Dung dịch chất X hòa tan \ce{Cu(OH)2}, thu được dung dịch màu xanh lam. Mặt khác, X bị thủy phân khi đun nóng trong môi trường acid. Chất X là
\choice
{glucose.}
{saccharose.}
{amylose.}
{fructose.}
\loigiai{}
\end{ex}
\begin{ex}
Ester X có công thức phân tử \ce{C4H8O2}. Thủy phân X trong dung dịch \ce{H2SO4} loãng, đun nóng, thu được sản phẩm gồm acetic acid và chất hữu cơ Y. Công thức của Y là
\choice
{\ce{HCOOH}.}
{\ce{C2H5OH}.}
{\ce{CH3COOH}.}
{\ce{CH3OH}.}
\loigiai{}
\end{ex}
\begin{ex}
Mùi tanh của cá là hỗn hợp các amine và một số tạp chất khác. Để khử mùi tanh của cá trước khi chế biến thực phẩm, nên áp dụng cách nào sau đây?
\choice
{Rửa cá bằng dung dịch nước muối.}
{Ngâm cá trong nước để amine tan vào nước.}
{Rửa cá bằng dung dịch soda.}
{Rửa cá bằng giấm ăn.}
\loigiai{}
\end{ex}
\begin{ex}
Khi thuỷ phân lipid trong môi trường base thì thu được muối của acid béo và
\choice
{ester đơn chức.}
{phenol.}
{alcohol đơn chức.}
{glycerol.}
\loigiai{}
\end{ex}
\begin{ex}
Giá trị pH mà khi đó amino acid có nồng độ ion lưỡng cực là cực đại được gọi là điểm đẳng điện (kí hiệu là pI). Khi $\text{pH} < \text{pI}$ thì amino acid đó tồn tại chủ yếu ở dạng cation, còn khi $\text{pH} > \text{pI}$ thì amino acid đó tồn tại chủ yếu ở dạng anion. Khi đặt trong một điện trường dạng anion sẽ di chuyển về cực $(+)$ còn dạng cation sẽ di chuyển về cực $(-)$. Cho các giá trị pI: glycine ($6{,}0$); glutamic acid ($3{,}2$); lysine ($9{,}7$). Giá trị pH tối ưu nhất để tách ba chất trên ra khỏi dung dịch hỗn hợp là
\choice
{$\text{pH} = 14{,}0$.}
{$\text{pH} = 9{,}7$.}
{$\text{pH} = 3{,}2$.}
{$\text{pH} = 6{,}0$.}
\loigiai{}
\end{ex}
\begin{ex}
Tripeptide là hợp chất mà phân tử có
\choice
{ba liên kết peptide, hai gốc $\alpha$-amino acid.}
{hai liên kết peptide, ba gốc $\alpha$-amino acid.}
{ba liên kết peptide, ba gốc $\alpha$-amino acid.}
{hai liên kết peptide, ba gốc $\beta$-amino acid.}
\loigiai{}
\end{ex}
\begin{ex}
Chất nào sau đây vừa tác dụng được với \ce{H2NCH2COOH}, vừa tác dụng được với \ce{CH3NH2}?
\choice
{\ce{NaOH}.}
{\ce{HCl}.}
{\ce{NaCl}.}
{\ce{CH3OH}.}
\loigiai{}
\end{ex}
\begin{ex}
Phần không phân cực (``đuôi'' kị nước) trong xà phòng và chất giặt rửa là
\choice
{nhóm sulfonate.}
{gốc hydrocarbon dài.}
{nhóm sulfate.}
{Nhóm carboxylate.}
\loigiai{}
\end{ex}
\begin{ex}
Glucose là một loại monosaccharide có nhiều trong quả nho chín. Công thức phân tử của glucose là
\choice
{\ce{C6H12O6}.}
{\ce{(C6H10O5)_n}.}
{\ce{C2H4O2}.}
{\ce{C12H22O11}.}
\loigiai{}
\end{ex}
\begin{ex}
Nhóm chức nào sau đây không có trong cấu tạo của glucose?
\choice
{Aldehyde.}
{Ketone.}
{Hemiacetal.}
{Hydroxy.}
\loigiai{}
\end{ex}
\begin{ex}
Ở động vật, tinh bột được dự trữ dưới dạng glycogen ở đâu?
\choice
{Ruột}
{Dạ dày}
{Gan}
{Máu}
\loigiai{}
\end{ex}
\begin{ex}
Số ester có cùng công thức phân tử \ce{C3H6O2} là
\choice
{4.}
{3.}
{2.}
{5.}
\loigiai{}
\end{ex}
\begin{ex}
Monosaccharide X được dùng trong công nghiệp để tráng bạc lên bề mặt thuỷ tinh trong sản xuất ruột phích. Cùng với Ag, sản phẩm hữu cơ được tạo thành khi cho X tác dụng với lượng dư dung dịch \ce{AgNO3} trong \ce{NH3} là
\choice
{ammonium gluconate.}
{ammonium carbonate.}
{gluconic acid.}
{khí carbon dioxide.}
\loigiai{}
\end{ex}
\begin{ex}
Chất có công thức phân tử \ce{C12H22O11}, được tạo thành trong quá trình thuỷ phân không hoàn toàn amylose có trong tinh bột là
\choice
{maltose.}
{glucose.}
{saccharose.}
{fructose.}
\loigiai{}
\end{ex}

\vspace{6pt}
\noindent{\color{mainGreen}\textbf{\large Phần II: Câu hỏi trả lời đúng sai.}} \\
\textit{Thí sinh trả lời từ câu 1 đến câu 2. Trong mỗi ý a), b), c), d) ở mỗi câu, thí sinh chọn đúng hoặc sai.}
\vspace{4pt}

\begin{ex}
Tiến hành thí nghiệm thủy phân ester theo các bước sau:\\
- Bước 1: Cho vào hai ống nghiệm mỗi ống 2 mL ethyl acetate.\\
- Bước 2: Thêm 2 mL dung dịch \ce{H2SO4} 20\% vào ống thứ nhất; 4 mL dung dịch \ce{NaOH} 30\% vào ống thứ hai.\\
- Bước 3: Lắc đều cả hai ống nghiệm, lắp ống sinh hàn, đun sôi nhẹ trong khoảng 5 phút, để nguội.
\choiceTF
Phản ứng xảy ra trong ống nghiệm 1 nhanh hơn phản ứng trong ống 2.
Vai trò của dung dịch \ce{H2SO4} 20\% và dung dịch \ce{NaOH} 30\% trong thí nghiệm trên là chất xúc tác và không bị tiêu hao sau phản ứng.
\True Sau bước 2, chất lỏng trong cả hai ống nghiệm đều phân thành hai lớp.
\True Sau bước 3: ống nghiệm (1) thể tích lớp chất lỏng phía trên giảm. Ống nghiệm (2) tạo thành hỗn hợp đồng nhất.
\loigiai{}
\end{ex}

\begin{ex}
Hình dưới đây mô tả cấu tạo phân tử saccharose gồm một gốc $\alpha$-glucose và một gốc $\beta$-fructose liên kết với nhau qua liên kết $\alpha$-1,2-glycoside:
\begin{center}
\includegraphics[width=0.65\textwidth]{C:/Users/Admin/.gemini/antigravity/brain/c29808c0-7d22-4aec-89d3-7a04648f0518/.user_uploaded/media_1794322695260_e258ffd0.jpg}
\end{center}
\choiceTF
Ngoài cấu tạo dạng mạch vòng, saccharose còn có cả cấu tạo ở dạng mở vòng của đơn vị cấu tạo glucose và fructose.
Mỗi phân tử saccharose gồm một đơn vị $\alpha$-glucose và một đơn vị $\alpha$-fructose liên kết với nhau qua liên kết $\alpha$-1,2-glycoside.
\True Saccharose và maltose đều có 42,1\% C về khối lượng.
Phân tử saccharose không có nhóm -OH hemiacetal, nhưng có nhóm -OH hemiketal.
\loigiai{}
\end{ex}

\vspace{6pt}
\noindent{\color{mainGreen}\textbf{\large Phần III: Câu trắc nghiệm trả lời ngắn.}} \\
\textit{Thí sinh trả lời từ câu 1 đến câu 4.}
\vspace{4pt}

\begin{ex}
Ethanol có thể được sản xuất từ cellulose hoặc tinh bột. Loại ethanol này được dùng để sản xuất xăng E10 (xăng chứa 10\% ethanol về thể tích). Lượng ethanol thu được từ 1 tấn mùn cưa (chứa 50\% cellulose, phần còn lại là chất trơ) có thể dùng để pha chế bao nhiêu L xăng E10? Biết hiệu suất quá trình sản xuất ethanol từ cellulose là 60\% và ethanol có khối lượng riêng là 0,789 g/mL. \textit{(Kết quả làm tròn đến hàng đơn vị)}.
\shortans{2159}
\loigiai{
\textbf{Trả lời: 2159}\\
$m_{\text{cellulose}} = 1000 \times 50\% = 500\text{ kg} = 500\,000\text{ gam}$.\\
Sơ đồ: $\ce{(C6H10O5)_n -> n C6H12O6 -> 2n C2H5OH}$\\
$m_{\text{ethanol}} = 500\,000 \times \frac{92}{162} \times 60\% = 170\,370{,}37\text{ gam}$.\\
$V_{\text{ethanol}} = \frac{170\,370{,}37}{0{,}789} = 215\,932\text{ mL} = 215{,}932\text{ L}$.\\
$V_{\text{xăng E10}} = \frac{215{,}932}{0{,}10} = 2159{,}32\text{ L} \approx \mathbf{2159}$.
}
\end{ex}

\begin{ex}
Cho các chất sau: ethyl acetate, tripalmitin, saccharose, ethylamine, Gly-Ala, albumin, Ala-Ala-Val. Số chất tham gia phản ứng thủy phân trong môi trường kiềm là bao nhiêu?
\shortans{5}
\loigiai{
\textbf{Trả lời: 5}\\
Các chất thủy phân trong môi trường kiềm gồm 5 chất:
\begin{enumerate}[noitemsep,topsep=1pt]
\item Ethyl acetate (ester).
\item Tripalmitin (triglyceride / chất béo).
\item Gly-Ala (dipeptide).
\item Albumin (protein).
\item Ala-Ala-Val (tripeptide).
\end{enumerate}
(Saccharose không bị thủy phân trong kiềm; ethylamine là amine không bị thủy phân).
}
\end{ex}

\begin{ex}
Cho các dung dịch sau: \ce{HCl}, \ce{NaOH}, \ce{CH3COOH}, \ce{CuSO4}. Có bao nhiêu dung dịch tác dụng được với ethylamine?
\shortans{3}
\loigiai{
\textbf{Trả lời: 3}\\
Ethylamine (\ce{C2H5NH2}) tác dụng với 3 dung dịch:
\begin{itemize}[noitemsep,topsep=1pt]
\item \ce{HCl}: $\ce{C2H5NH2 + HCl -> C2H5NH3Cl}$.
\item \ce{CH3COOH}: $\ce{C2H5NH2 + CH3COOH -> CH3COONH3C2H5}$.
\item \ce{CuSO4}: tạo kết tủa \ce{Cu(OH)2} rồi hòa tan tạo phức xanh lam $[\ce{Cu(C2H5NH2)4}]^{2+}$.
\end{itemize}
(\ce{NaOH} không phản ứng vì cùng có tính base).
}
\end{ex}

\begin{ex}
Một phân tử saccharose có bao nhiêu nguyên tử hydrogen?
\shortans{22}
\loigiai{
\textbf{Trả lời: 22}\\
Công thức phân tử saccharose là \ce{C12H22O11}, do đó có 22 nguyên tử hydrogen.
}
\end{ex}

\vspace{6pt}
\noindent{\color{mainGreen}\textbf{\Large B. PHẦN TỰ LUẬN}}
\vspace{4pt}

\noindent\textbf{Câu 1.} Hoàn thành phương trình phản ứng sau:\\
\ce{CH3COOC6H5} (phenyl acetate) + \ce{NaOH} dư (dung dịch, đun nóng).

\vspace{6pt}
\multido{\i=1+1}{4}{\noindent\dotfill\\[4pt]}

\vspace{6pt}
\noindent\textbf{Câu 2.} Glucose tồn tại đồng thời dạng mạch hở và mạch vòng ($\alpha$ và $\beta$). Glucose có tính chất của aldehyde và của polyalcohol. Cho các chất sau: \ce{Cu(OH)2}, nước bromine, dung dịch \ce{I2} trong \ce{KI}, thuốc thử Tollens và \ce{CH3OH}/\ce{HCl} khan. Ở điều kiện thích hợp, liệt kê các chất tác dụng được với glucose?

\vspace{6pt}
\multido{\i=1+1}{5}{\noindent\dotfill\\[4pt]}

\vspace{6pt}
\noindent\textbf{Câu 3.} Viết công thức cấu tạo của Alanine.

\vspace{6pt}
\multido{\i=1+1}{3}{\noindent\dotfill\\[4pt]}

\vspace{6pt}
\noindent\textbf{Câu 4.} Trong cơ thể, glucose là sản phẩm của quá trình thủy phân tinh bột dưới tác dụng của enzyme. Trong quá trình hô hấp tế bào, 1 mol glucose bị oxi hóa hoàn toàn qua các phản ứng sinh hóa tạo thành carbon dioxide, nước và giải phóng năng lượng là 2880 kJ. Trong khẩu phần ăn hằng ngày của một người trưởng thành cần 300 gam gạo. Hàm lượng tinh bột trong gạo là 80\%. Giả sử 60\% lượng tinh bột trong gạo chuyển hóa thành glucose và toàn bộ lượng glucose này tham gia vào quá trình hô hấp tế bào. Năng lượng giải phóng từ sự oxi hóa hoàn toàn glucose chiếm bao nhiêu \% so với nhu cầu năng lượng trung bình (2000 Calori/ngày) của một người trưởng thành? \textit{(Kết quả làm tròn đến hàng phần mười)}.\\
\textit{Cho biết: 1 Calori = 1 kcal = 1000 cal = 4184 J.}

\vspace{6pt}
\multido{\i=1+1}{6}{\noindent\dotfill\\[4pt]}

\vspace{6pt}
\noindent\textbf{Câu 5.} Amine X có phân tử khối nhỏ hơn 80. Trong phân tử X, nitrogen chiếm 19,18\% về khối lượng. Cho X tác dụng với nitrous acid (\ce{HNO2}) ở nhiệt độ thường thu được alcohol Y. Oxi hóa không hoàn toàn Y thu được ketone Z. Viết công thức cấu tạo của Y.

\vspace{6pt}
\multido{\i=1+1}{6}{\noindent\dotfill\\[4pt]}

\vspace{6pt}
\noindent\textbf{Câu 6.} Tiến hành thí nghiệm với các chất X, Y, Z là một trong các chất sau: hồ tinh bột, albumin, ethyl acetate, vinyl acetate. Kết quả được ghi ở bảng sau:
\begin{center}
\renewcommand{\arraystretch}{1.2}
\begin{tabular}{|c|p{8.5cm}|c|}
\hline
\rowcolor{tableGreen}
\textbf{Mẫu thử} & \textbf{Thí nghiệm} & \textbf{Hiện tượng} \\
\hline
\textbf{X} & Tác dụng với \ce{Cu(OH)2} trong môi trường kiềm & Có màu tím \\
\hline
\textbf{Y} & Đun nóng với dung dịch \ce{NaOH} loãng (vừa đủ). Thêm tiếp dung dịch \ce{AgNO3} trong \ce{NH3}, đun nóng & Tạo kết tủa Ag \\
\hline
\textbf{Z} & Tác dụng với dung dịch \ce{I2} loãng & Có màu xanh tím \\
\hline
\end{tabular}
\end{center}
Tìm các chất X, Y, Z.

\vspace{6pt}
\multido{\i=1+1}{5}{\noindent\dotfill\\[4pt]}

\begin{center}
\textbf{--- HẾT ---}
\end{center}
\end{document}
